\chapter*{Abstract}

This thesis presents an advanced architectural framework for Conversational AI systems leveraging the \textbf{Retrieval-Augmented Generation (RAG)} paradigm combined with distributed graph databases. The primary objective is to mitigate common LLM limitations, such as ``hallucinations'' and knowledge cutoff, by providing a verifiable external ``source of truth''.

The proposed system utilizes a two-tier data management strategy: a \textbf{Global Remote Storage} for long-term knowledge and a \textbf{Local Shard} acting as a high-performance cache to minimize retrieval latency. By employing \textbf{Neo4j} and hybrid vector-graph search techniques, the system enables efficient navigation of complex relationships within medical datasets. The integration includes a sophisticated synchronization logic driven by parameterized Cypher queries and a Least Frequently Used (LFU) eviction policy. The final prototype, built with \textbf{Unreal Engine} and the \textbf{FANTASIA} framework, incorporates \textbf{Whisper} for ASR and \textbf{ElevenLabs} for TTS, creating a seamless voice-driven interaction. Experimental analysis confirms that the distributed approach yields a substantial reduction in end-to-end latency, validating its effectiveness for real-time applications.

\chapter*{Sommario}

Il presente lavoro di tesi propone un’architettura innovativa per sistemi di Conversational AI basata sul paradigma \textbf{Retrieval-Augmented Generation (RAG)}, integrando database a grafo distribuiti per superare i limiti di memoria e affidabilità dei modelli linguistici di grandi dimensioni (LLM). Il cuore del progetto risiede in una gestione distribuita dei dati strutturata su due livelli: un \textbf{Remote Storage} globale e uno \textbf{Shard Locale} che funge da cache ad alte prestazioni.

Attraverso l’utilizzo di \textbf{Neo4j} per la persistenza del grafo e di modelli come \textbf{Llama 3} e \textbf{Gemma} per il ragionamento e gli embedding, il sistema permette una navigazione semantica e contestuale avanzata. L’implementazione include logiche di sincronizzazione dinamica e algoritmi di \textit{eviction} LFU per mantenere l’efficienza dello shard. L’interfaccia utente, sviluppata in \textbf{Unreal Engine}, integra pipeline vocali con \textbf{Whisper} ed \textbf{ElevenLabs}, garantendo un'interazione naturale. I risultati sperimentali dimostrano che l'approccio distribuito riduce significativamente la latenza delle risposte, migliorando l'esperienza utente complessiva.